% Generated by `taskgen course examples` from modular.fermat_dim1 (seed 5).
% Do not edit: change examples.yaml and regenerate.
\begin{example}[код с уменьшением на единицу]\label{ex:dim1-code}
В кольце $\mathbb{Z}/257\mathbb{Z}$, $257 = 2^{8} + 1$, вычеты хранятся в коде с уменьшением на единицу ($9$ бит: старший бит --- признак нуля, младшие $8$ бит --- $x - 1$ при $x \ne 0$). Даны коды \[a = 011011000_{2f}, \qquad b = 001110011_{2f}.\] Вычислить коды суммы $a + b$ и произведения $a \cdot b$, выполняя операции над кодами. Проверить ответ, перейдя к обычной записи.

\par\noindent\textit{Решение.} Обозначим $x' = x - 1$ ($8$-битовую часть кода ненулевого $x$). Тогда $(x + y)' \equiv x' + y' + 1$, $(xy)' \equiv xy - 1 = x'y' + x' + y'$ (по модулю $257$), и $2^{8} \equiv -1$.\par Сумма: $a' + b' = 11011000_2 + 1110011_2 = 101001011_2$ ($216 + 115 = 331$). Есть перенос из $8$-го разряда: $2^{8} \equiv -1$ съедает прибавляемую единицу, код суммы $= a' + b' - 2^{8} = 75$.\par Произведение: $P = a'b' + a' + b' = 24840 + 216 + 115 = 25171 = 110001001010011_2$, $P = H \cdot 2^{8} + L$, $H = 98$, $L = 83$; так как $2^{8} \equiv -1$, код произведения $\equiv L - H = -15$, прибавляем $257$: $242$.\par Проверка: $a = 217$, $b = 116$, $a + b \equiv 76$, $a b = 25172 \equiv 243 \pmod{257}$.\[a + b = 001001011_{2f}, \qquad a \cdot b = 011110010_{2f}.\]
\end{example}
